\documentclass[10pt,a4paper]{amsart}
\usepackage[margin=19mm]{geometry}
\usepackage{amsmath,amssymb,mathtools}
\usepackage[scr=rsfs,cal=euler]{mathalfa}
\usepackage{xfrac}
\usepackage[unicode,hidelinks,bookmarks=false]{hyperref}
\newcommand{\ie}{i.e.\ }
\newcommand{\cf}{cf.\ }
\newcommand{\resp}{respectively\ }
\newcommand{\Subsection}{\subsection}
\providecommand{\nobreakdash}{\nobreak\hbox{-}}
\setlength{\emergencystretch}{2em}
\multlinegap0pt
\usepackage{bm}
\usepackage{bbm}
\usepackage{dsfont}
\usepackage{xspace}
\usepackage{cancel}
\usepackage{mathtools}
\usepackage{amsthm}
\makeatletter
\@ifundefined{assumption}{\newtheorem{assumption}{Assumption}}{}
\@ifundefined{lemma}{\newtheorem{lemma}{Lemma}}{}
\@ifundefined{theorem}{\newtheorem{theorem}{Theorem}}{}
\makeatother
\title[Towards regularized learning from functional data with covar]{Towards regularized learning from functional data with covariate shift}
\author{Markus Holzleitner, Sergiy Pereverzyev Jr., Sergei V. Pereverzyev, Vaibhav Silmana, S. Sivananthan}
\date{Source: arXiv:2601.21019v1 (CC BY 4.0)}
\begin{document}
\maketitle

\noindent\emph{Source and licence.} Towards regularized learning from functional data with covariate shift. Version arXiv:2601.21019v1 (CC BY 4.0), distributed under the Creative Commons Attribution 4.0 International licence, as recorded in the arXiv licence metadata for this exact version and shown on the abstract page https://arxiv.org/abs/2601.21019v1. Full rights notice: licenses/arxiv-2601.21019-CC-BY-4.0.txt.\par\smallskip

\noindent\textit{Editorial scope.} Counted unit: the Theorem whose source label is \texttt{vrkhs:theorem:proof} in the article (source0.tex lines 1148--1178, source label \texttt{vrkhs:theorem:proof}), delivered together with the article's own complete original proof of it (source0.tex lines 1180--1283, one \texttt{proof} environment of 104 lines). No mathematics is written, repaired or re-derived: statement and proof are the article's own text line for line. Required context retained uncounted: kernel-bound (lines 806--811); assump:source (lines 853--862); eq:condition\_on\_lambda (lines 961--964); L\_2(lemma-A) (lines 977--1010); L2:bound-B (lines 1012--1025). Every definition, display and label of those lines is retained unchanged. Omitted from this package and not redistributed: 1--805, 812--852, 863--960, 965--976, 1011--1011, 1026--1147, 1179--1179, 1284--2676. Nothing inside a retained excerpt is omitted or altered: every retained body is delivered exactly as the article's lines are written, so no omission receipt of any kind accompanies this package. Citations: the retained excerpts carry no citation command at all, so no bibliography entry of the article is transcribed and no reference is redistributed. Reference adaptation: none was needed and none was made; every reference inside the delivered unit resolves inside it, so no unresolved reference or citation is produced. Printed slips kept literal: no printed word, symbol, hypothesis or number of the counted statement or of its proof was altered. Layout and class: the article's own class is \texttt{article} with packages including jmlr2e, lastpage, booktabs, amsmath; the wrapper is the lane's pinned \texttt{amsart} editorial wrapper at 10pt on A4, which restates the theorem-like environments the retained text needs and adds the article's own macro definitions that the retained text needs (none), with extra wrapper packages (bm, bbm, dsfont, xspace, cancel, mathtools). The delivered numbering is the unit's own. Assets: no retained span contains a figure environment, a graphics-inclusion command, a PDF-page inclusion, a table or a TikZ picture; no image file is copied into this package, the package declares no asset dependency and no image file is redistributed with it. Licence: version arXiv:2601.21019v1 of this article is distributed under the Creative Commons Attribution 4.0 International licence, as recorded in the arXiv licence metadata for this exact version and shown on the abstract page https://arxiv.org/abs/2601.21019v1. A full rights notice is delivered at licenses/arxiv-2601.21019-CC-BY-4.0.txt. Characters: every retained excerpt is pure ASCII, so the delivered engine, XeLaTeX with its default Latin Modern text font, reports no missing character.
\begin{assumption}\label{kernel-bound}
Let $k : X \times X \to \mathbb{R}$ be the reproducing kernel of a scalar-valued RKHS $\mathcal{H}$, satisfying
\[
k(x,x) \le \kappa^2 \quad \mathrm{for\ all}\ x \in X .
\]
\end{assumption}

\begin{assumption}{(General source condition)}\label{assump:source}
{
We assume that the operator $C_\ast$ belongs to the Hilbert--Schmidt class
$\mathcal S_2(\mathcal H^{\varphi},Y)$, where $\varphi$ is the index function and  the space $\mathcal H^{\varphi}$ is defined as
\[
\mathcal H^{\varphi}
:= \left\{ f \in \mathcal H \;:\; f = \varphi(C_X) g \quad \mathrm{for \, some } \; g \in \mathcal H \right\}.
\]
}
\end{assumption}

\begin{equation}\label{eq:condition_on_lambda}
\mathcal{B}_{n,\lambda}
\log\!\left(\frac{2}{\delta}\right)\leq \frac{\sqrt{\lambda}}{2}.
\end{equation}

\begin{lemma}\label{L_2(lemma-A)}
Suppose that {assumptions~\ref{kernel-bound}--\ref{assump:source}} are satisfied.
If  the qualification $\nu$ of the regularization family covers the index function \(\varphi(t)\sqrt{t}\) and \eqref{eq:condition_on_lambda} holds,
then with probability at least \(1-\delta\), we have

\[\big\|(C_*-\bar{C}_{\lambda}^\beta)\, C_X^{1/2} \big\|_{\mathcal S_2(\mathcal H,Y)} \leq\;
\left(
\frac{A_1}{n}\sqrt{\frac{1}{\lambda}}
+ A_2\sqrt{\frac{\mathcal{N}(\lambda)}{n}}
\right) \log\left(\frac{2}{\delta}\right)
+ A_3\,\sqrt{\lambda}\varphi(\lambda).\]

Here $\bar{C}_{\lambda}^{\beta}
= C_* \widehat{C}_X^{\beta} g_{\lambda}(\widehat{C}_X^{\beta})$,
and the constants $A_1, A_2, A_3$ are given by
\begin{align*}
A_1 &:= 2^{\,2\nu-\bar{\nu}+2}(\gamma_\nu+\gamma)\,
      \max\!\left(1,\frac{1}{c}\right)\,
      \bar{\nu}\, b\, \kappa^2\,
      \| C_{*}\|_{\mathcal S_2(\mathcal H,Y)},\\[4pt]
A_2 &:= 2^{\,2\nu-\bar{\nu}+1}(\gamma_\nu+\gamma)\,
      \max\!\left(1,\frac{1}{c}\right)\,
      \bar{\nu}\,\sqrt{b}\,\kappa\,
      \|C_{*}\|_{\mathcal S_2(\mathcal H,Y)},\\[4pt]
A_3 &:= 2^{\,2\nu-\bar{\nu}}(\gamma_\nu+\gamma)\,
      \max\!\left(1,\frac{1}{c}\right)\,
      \|C_{o}\|_{\mathcal S_2(\mathcal H,Y)},\\[4pt]
\end{align*}
where $\bar{\nu} := \lfloor \nu -1/2 \rfloor$ denotes the greatest integer
less than or equal to~$\nu -1/2$.


    
\end{lemma}

\begin{lemma}\label{L2:bound-B}
Suppose that assumptions~\ref{kernel-bound}--\ref{assump:source} are satisfied and that equation~\eqref{eq:condition_on_lambda} holds. Then, with probability at least \(1-\delta\), we have
\begin{align*}
\bigl\| ( \widehat{C}_\lambda^\beta-\bar{C}_{\lambda}^\beta)\, C_X^{1/2} \bigr\|_{\mathcal S_2(\mathcal H,Y)}
\;&\leq\;
(B+D)\left(
\frac{10bM\kappa}{n\sqrt{\lambda}}
+ 5
M\sqrt{
\frac{b \; \mathcal{N}(\lambda)}{n}
}
\right)\log\!\left( \frac{2}{\delta} \right).
\end{align*}
\end{lemma}

\begin{theorem}\label{vrkhs:theorem:proof}
Suppose that {assumptions~\ref{kernel-bound}--\ref{assump:source}} are satisfied.
If \eqref{eq:condition_on_lambda} holds and the qualification $\nu$ of the regularization family covers the index function \(\varphi(t)\),
then with probability at least \(1-\delta\), we have

\[
\big\|f^*-\widehat{f}_{\textbf{z},\lambda}\big\|_{\mathcal G}\leq \left(
\frac{A_1}{n\lambda}
+ A_2\sqrt{\frac{\mathcal{N}(\lambda)}{n\lambda}}
\right) \log\left(\frac{2}{\delta}\right)
+ A_3\,\varphi(\lambda).
\]

\noindent
Here the constants $A_1,A_2,A_3$ are given by 
\begin{align*}
A_1
&= 2^{2\nu- \bar \nu+2}(\gamma_\nu+\gamma)\max\!\left(1,\frac{1}{c}\right)
\,\bar \nu\, b\,\kappa^{2}\,\|C_{*}\|_{\mathcal S_2(\mathcal H,Y)}
\;+\; 2\sqrt{10}\,b\,M\,\kappa\,(B+D), \\[6pt]
A_2
&= 2^{2\nu- \bar \nu+1}(\gamma_\nu+\gamma)\max\!\left(1,\frac{1}{c}\right)
\, \bar \nu\, \sqrt{b}\kappa\,\|C_{*}\|_{\mathcal S_2(\mathcal H,Y)}
\;+\; \sqrt{10}\,\sqrt{b}\,M(B+D), \\[6pt]
A_3
&= 2^{2\nu- \bar \nu}(\gamma_\nu+\gamma)\max\!\left(1,\frac{1}{c}\right)
\,\|C_{o}\|_{\mathcal S_2(\mathcal H,Y)}.
\end{align*}
where $\bar{\nu} := \lfloor \nu  \rfloor$ denotes the greatest integer
less than or equal to~$\nu $.
\end{theorem}

\begin{proof}
    \noindent
We have
\begin{align*}
\big\|f^*-\widehat{f}_{\textbf{z},\lambda}\big\|_{\mathcal G}
&=
\big\|C_*-\widehat{C}_{\lambda}^\beta\, \big\|_{\mathcal S_2(\mathcal H,Y)} \leq
\underbrace{\big\|C_*-\bar{C}_{\lambda}^\beta\, \big\|_{\mathcal S_2(\mathcal H,Y)}}_{\textbf{A}}  +
\underbrace{\big\|\bar{C}_{\lambda}^\beta -\widehat{C}_\lambda^\beta\, \big\|_{\mathcal S_2(\mathcal H,Y)}}_{\textbf{B}}. 
\end{align*}

\noindent
The bound for Term~\textbf{A} can be obtained by following the steps of
Lemma~\ref{L_2(lemma-A)}. Let $\bar \nu$ denote the greatest integer less than $\nu$.
Then, with probability at least $1-\delta$, we obtain

\[
\begin{aligned}
\big\|C_*-\bar{C}_{\lambda}^\beta\big\|_{\mathcal  S_2(\mathcal H,Y)}
\leq\;&\max\!\left(1,\frac{1}{c}\right)
\Bigg(
\frac{2^{\,2\nu-\bar \nu}(\gamma_\nu+\gamma)\,
      4\,\bar\nu\,b\,\kappa^2\,
      \|C_{*}\|_{\mathcal S_2(\mathcal H,Y)}}{n\lambda}
\\
&\quad
+ 2^{\,2\nu-\bar\nu}(\gamma_\nu+\gamma)\,
      2\,\bar\nu\,\sqrt{b}\kappa\,
      \|C_{*}\|_{\mathcal S_2(\mathcal H,Y)}
      \sqrt{\frac{\mathcal{N}(\lambda)}{n\lambda}}
\Bigg)
\log\!\left(\frac{2}{\delta}\right)
\\
&\quad
+ 2^{\,2\nu-\bar\nu}(\gamma_\nu+\gamma)\,
      \max\!\left(1,\frac{1}{c}\right)\,
      \|C_{o}\|_{\mathcal S_2(\mathcal H,Y)}\,
      \varphi(\lambda).
\end{aligned}
\]



To obtain bound for $Term~$\textbf{B}, we observe
\begin{align*}
\bigl\|  \widehat{C}_\lambda^\beta-\bar{C}_{\lambda}^\beta\,  \bigr\|_{\mathcal  S_2(\mathcal H,Y)}
&=
\bigl\|
( \widehat{C}_{Y,X}^\beta-{C_*}\widehat C_X^\beta)\,
g_\lambda(\widehat{C}_X^\beta)\,
\bigr\|_{\mathcal  S_2(\mathcal H,Y)} .
\end{align*}

\noindent
Next, we split the operator as
\begin{align*}
&\bigl\|
( \widehat{C}_{Y,X}^\beta-{C_*}\widehat C_X^\beta)\,
g_\lambda(\widehat{C}_X^\beta)\,
\bigr\|_{\mathcal S_2(\mathcal H,Y)} \\
&\leq
\underbrace{\bigl\|
( \widehat{C}_{Y,X}^\beta-{C_*}\widehat C_X^\beta)\,
(C_{X}+\lambda I)^{-1/2}
\bigr\|_{\mathcal S_2(\mathcal H,Y)}}_{A}
\\
&\hspace{9.5em}\times
\underbrace{\bigl\|
(C_{X}+\lambda I)^{1/2}
(\widehat{C}^{\beta}_{X}+\lambda I)^{-1/2}
\bigr\|_{op}}_{B}
\;
\underbrace{\bigl\|
(\widehat{C}^{\beta}_{X}+\lambda I)^{1/2}
g_\lambda(\widehat{C}_X^\beta)
\bigr\|_{op}}_{C}.
\end{align*}

\noindent
The bound for $Term~A,B$ are already given in proof of Lemma \ref{L2:bound-B}. For $Term~C$, we can write

\[
C = {sup}_\sigma \; | (\sigma +\lambda)^{1/2} g_\lambda(\sigma)| \leq {sup}_\sigma \; \left| \frac{ (\sigma +\lambda) g_\lambda(\sigma)}{(\sigma +\lambda)^{1/2}}\right| \leq \frac{B+D}{\sqrt{\lambda}}.   \]

\noindent
Hence, with probability at least $1-\delta$, we obtain

\begin{align*}
\bigl\|  \widehat{C}_\lambda^\beta-\bar{C}_{\lambda}^\beta\,  \bigr\|_{\mathcal  S_2(\mathcal H,Y)}
\;&\leq\;
\sqrt{\frac{5}{2}}(B+D)\left(
\frac{4bM\kappa}{n{\lambda}}
+ 2
M\sqrt{
\frac{b \; \mathcal{N}(\lambda)}{n\lambda}
}
\right)\log\!\left( \frac{2}{\delta} \right)
.
\end{align*}

\noindent
Finally, combining the bounds for Term~\textbf{A} and Term~\textbf{B} yields the desired result.

\end{proof}

\end{document}
